\section{Integral cusp lattices on the original geometric family}
\label{sec:icl-v46}

This chapter concerns the lattice defined by \emph{all} the original
finite-place integral-frame tests.  Its first part proves finite generation
by constructing separating evaluations on the actual moduli space.  Its
second part isolates the analytic and integral-theta inputs needed to prove
that this lattice spans, and proves their exact consequence for
\texttt{LatFull}.  The distinction matters: the latter inputs are original
geometric producer statements, and are not consequences of a dimension
count or of the finite rational-theta implications in RTG.

Fix the source data \(D\), \(T:\mathrm{Core}(D)\), its original
\(Q,X_j,c,l\), a row \(j\in\mathrm{InR}(D)\), and a dominant weight
\(w\).  Put
\[
 K=c.F,\qquad E=D.\mathrm{host}.E,\qquad
 C_w=T.\mathrm{cuspA}(j,\mathrm{wtA}(w),w.n),\qquad
 \mathcal R=T.R(j).
\]
All scalar extensions below use the given inclusion of \(K\).  Auxiliary
number fields used to evaluate a section do not replace \(c\), the common
field \(D.\mathrm{what}_2.LL\), its fixed complex embedding, the chosen
Hermitian forms, the metric, the measure, or any tuple phase.  In particular
the last transport uses the original \(\mathrm{FWTie}\) and its map.

We use the following classical foundations with their stated scope.
The source interfaces and the exact places where they enter are recorded
in the accompanying foundation manifest.

\begin{enumerate}
\item Riemann's algebraic uniformization of polarized complex tori, the
algebraicity of complex abelian varieties with complex multiplication,
and their potential good reduction at every finite place.  For the last
assertion a primary reference is Serre--Tate, \emph{Good reduction of
abelian varieties}, Annals of Mathematics 88 (1968), Theorem 6(a).
\item The extension property for homomorphisms of abelian schemes over
a normal Dedekind base, the theorem of the cube for rigidified line
bundles, and cohomology and base change for an ample principal line
bundle on an abelian scheme.  These are used for the given polarization
and line bundle, including at primes ramified in \(E\).
\item The finite-dimensionality of holomorphic cuspidal automorphic
forms of a fixed highest weight on a fixed arithmetic quotient.  It is
used together with the original algebraic uniformization, not with a
newly selected space of automorphic functions.
\end{enumerate}

The theta-kernel and ramified jet inputs introduced later are kept
separate from these classical foundations.  In particular, this chapter
does not relabel the original kernel identity, chart bounds, or short
series gluing as published theorems.

% BEGIN ICL1
\paragraph{ICL1. The literal valuation lattice.}
\label{claim:ICL1}
For a finite prime \(\mathfrak p\) of \(K\), the source set
\(\mathrm{valSet}(X_j,c,\mathrm{algCls})(\mathfrak p)\) consists of the maps
\[
 f\longmapsto e^{-1}\operatorname{ord}_{V}
                 (\operatorname{ev}_{I}(y,f)),                 \tag{ICL.1}
\]
with value \(+\infty\) at zero.  Here \(V\) is a complete discrete
valuation ring, its residue field is finite, \(\iota:K\hookrightarrow
\operatorname{Frac}(V)\) has ramification index \(e>0\) over
\(\mathfrak p\), and \(y\) is an original \(\mathrm{IntFr}\).
An \(\mathrm{IntFr}\) includes an algebraic test presentation, an integral
semiabelian model, bases of each integral Lie eigenspace, and a generator
of its integral theta-section line with the source rigidified
identifications.  Thus the source \(\mathrm{supLat}(T,j,w)\) is exactly
\[
 L_w=\{f\in C_w:\ v(f)\geq0\text{ for every }\mathfrak p
           \text{ and every }v\in\mathrm{valSet}(\mathfrak p)\}. \tag{ICL.2}
\]
It is an \(\mathcal O_K\)-module.  Its finite generation and the equality
\(K L_w=C_w\) are precisely \(\mathrm{FullAt}(T,j,w)\).
\begin{proof}
The definition \(\mathrm{IntFr.IsAlg}\) is the finite-residue-field
condition; it imposes no fixed field of definition on the test point.
The definition of \(\mathrm{PtV}\) supplies \(e,\iota,y\), and
\(\mathrm{vAtV}\) is (ICL.1).  Conversely every such tuple is a
\(\mathrm{PtV}\) with its unique auxiliary \(\mathrm{Unit}\)-label,
so it occurs in \(\mathrm{valSet}\).  Evaluation is additive,
\(K\)-linear through \(\iota\), and multiplicative.  Consequently
\(v(f+g)\geq\min(v(f),v(g))\),
\(v(af)=\operatorname{ord}_{\mathfrak p}(a)+v(f)\) for \(a\ne0\),
and \(v(fg)=v(f)+v(g)\).  Since the prime valuations of an element
of \(\mathcal O_K\) are nonnegative, (ICL.2) is a module.  The
source \(\mathrm{latIn}\) is the image of the corresponding submodule
of \(C_w\) under its inclusion in \(\mathcal R\); that inclusion is
injective.  This proves both the set identification and the last
assertion, without assuming that the module already spans.
\end{proof}
% END ICL1

% BEGIN ICL2
\paragraph{ICL2. Finitely many integral evaluations suffice.}
\label{claim:ICL2}
Let \(C\) be a finite-dimensional \(K\)-vector space and
\(L\subseteq C\) an \(\mathcal O_K\)-submodule.  Suppose that for
every nonzero \(f\in C\) there are a finite number-field extension
\(K_f/K\) and a \(K\)-linear map \(\varphi_f:C\to K_f\) such that
\[
 \varphi_f(f)\ne0,\qquad \varphi_f(L)\subseteq\mathcal O_{K_f}.
                                                               \tag{ICL.3}
\]
Then \(L\) is finitely generated.  Neither saturation nor a spanning
assumption on \(L\) is needed.
\begin{proof}
Start with \(W_0=C\).  If \(W_i\ne0\), choose a nonzero
\(f_i\in W_i\), choose (ICL.3), and set
\(W_{i+1}=W_i\cap\ker\varphi_{f_i}\).  Its dimension is strictly
smaller.  In at most \(\dim_K C\) steps the common kernel is zero.
The resulting diagonal map embeds \(L\) into
\(\prod_i\mathcal O_{K_{f_i}}\).  Each factor is a finite
\(\mathcal O_K\)-module, and \(\mathcal O_K\) is Noetherian, so
its submodule \(L\) is finite.  If \(C=0\), use the empty product
and \(L=0\).  The argument concerns the whole module \(L\), not a
chosen finite-index submodule.
\end{proof}
% END ICL2

% BEGIN ICL3
\paragraph{ICL3. Rational Hermitian flags detect a nonzero section.}
\label{claim:ICL3}
Let \(X\) be one original Hermitian uniformization component, with
Hermitian space \(E^r\), CM type \(\Psi\), and signature
\((p_0,r-p_0)\) at every \(\tau\in\Psi\).  Rational positive
\(p_0\)-planes and rational orthogonal complementary negative planes
are dense in its period domain.  More precisely, if a continuous
section \(P\) on the original frame domain is Borel equivariant
with a nonzero character and is nonzero somewhere, there is a frame
\(x\) with \(P(x)\ne0\) whose positive and negative subspaces
come from orthogonal \(E\)-subspaces.  The spanning vectors can all
be taken in any prescribed full \(\mathcal O_E\)-lattice.
\begin{proof}
Choose a frame where \(P\ne0\).  Hermitian Gram--Schmidt, in the
order with the positive columns first, changes it by an invertible
Borel matrix.  The defining nonzero principal minors ensure that
each division is legal, and Borel equivariance preserves
nonvanishing.  It is therefore enough to approximate a frame with
orthogonal columns and the required signs.

The real-linear Minkowski map
\(E\otimes_{\mathbb Q}\mathbb R\to\prod_{\tau\in\Psi}\mathbb C\)
is an isomorphism.  Its image of \(E\) is dense: in a
\(\mathbb Q\)-basis this is the density of \(\mathbb Q^d\) in
\(\mathbb R^d\).  Approximate all columns at all embeddings by a
single matrix over \(E\).  Apply Gram--Schmidt over \(E\), using
the original Hermitian form.  The coefficients remain in \(E\),
since the scalar products and nonzero self-pairings lie there.
At all the finitely many embeddings this operation is continuous
near the chosen orthogonal frame.  A sufficiently close
approximation is still a basis, has the same signs, and lies in
the open nonvanishing set of \(P\).

Multiplying the finitely many resulting vectors by a common
positive integer puts them in the prescribed full lattice.  Keep
the approximating complex frame before this multiplication; its
subspaces are unchanged.  If the source
\(\mathrm{LamBasis}(X,\mathcal L,x,\lambda)\) notation is desired,
choose \(\mathcal L>0\) larger than the square roots of all the
finitely many absolute self-pairings.  All its qualitative
conditions hold.  No quantitative bound for \(\mathcal L\), no
maximum of \(P\), and no uniform lattice-basis radius is asserted.
The cases \(p_0=0\) and \(p_0=r\) use the same construction with
one of the two subspaces zero.
\end{proof}
% END ICL3

% BEGIN ICL4
\paragraph{ICL4. These flags give actual algebraic CM points.}
\label{claim:ICL4}
Under the original moduli uniformization, a frame supplied by
ICL3 lies over a CM point of the same \(X_j\).  Its underlying
polarized abelian variety, \(\mathcal O_E\)-action, level marking,
rigidified principal line bundle and any chosen algebraic frame
are defined over some finite extension of \(K\).  Every nonzero
original section has a nonzero restriction to the frame fibre
over at least one such point.
\begin{proof}
The rational orthogonal decomposition gives \(r\) one-dimensional
\(E\)-Hodge subspaces.  The lattice of the uniformized torus is
commensurable with the sum of its intersections with these
subspaces: clear denominators in the finitely many rational
projection matrices.  Thus the torus is isogenous to a product
of the corresponding CM tori.  Equivalently the product algebra
\(E^r\), of dimension twice the abelian dimension over
\(\mathbb Q\), acts on its rational Hodge structure.  Riemann
uniformization turns these rational Hodge endomorphisms into
algebraic endomorphisms.  The variety is therefore of CM type.
CM algebraicity gives a model over a number field.  Enlarge that
field to contain the fixed copy of \(K\) and to define the
finitely many action maps, polarization, line-bundle data and
level-torsion coordinates.  These are algebraic data of finite
presentation.  Their relations hold after complex base change
and hence over this field.  The moduli point is in the original
component with the original marking; an isogenous object is
used to prove its CM property, not substituted for it.

A natural algebraic rule restricted to a fixed frame fibre is a
regular function on
\(\prod_\tau\mathrm{GL}_{m_\tau}\times\mathbb G_m\).
Indeed, apply the rule to the universal matrices and the
universal invertible section scalar.  Their coordinate ring is
a finitely generated reduced \(K\)-algebra, and therefore is
one of the source test rings after a presentation.  Naturality
identifies every later evaluation with evaluation of this one
regular function.  This also proves compatibility with the
source \(\mathrm{FrChg}((Y_\tau)^t,s)\) convention.

Finally, a nonzero rule in the original section space has a
nonzero complex reading.  To check this assertion directly,
its values on a reduced finite-type test algebra are detected
by its complex points; uniformization and the nonzero Borel
weight express each such value as a scalar times a frame
reading.  Vanishing of all readings would make the rule zero
on every test algebra.  Apply ICL3 to a nonzero reading.  At its
CM point the regular function on the frame fibre is nonzero,
since its value at that complex frame is nonzero.
\end{proof}
% END ICL4

% BEGIN ICL5
\paragraph{ICL5. One integral model after a finite extension.}
\label{claim:ICL5}
Let \(x\) be a CM point from ICL4.  After a finite extension
\(K'/K\), the actual abelian variety over \(x\) extends to an
abelian scheme over \(\mathcal O_{K'}\), with its original action
and rigidified symmetric principal line bundle.  Its theta
section module and each integral Lie eigenspace are finite
projective modules of the expected ranks.  They commute with
further flat extensions of the base.
\begin{proof}
There are only finitely many places of bad reduction.  Potential
good reduction for a CM abelian variety gives a finite local
extension at each.  A finite global extension realizing these
local extensions can be chosen by weak approximation and
Krasner's lemma.  Take a normal closure to handle all places
above each chosen prime.  Good reduction is preserved by field
extension, so the resulting variety has good reduction
everywhere.  Its local abelian models glue by the extension
and uniqueness property of abelian schemes.  The finitely
many original action maps and their relations extend by the
same property.

For clarity, no change of the generic theta line is needed.
On the regular abelian scheme, close a divisor representing
that line on the generic fibre.  Regularity makes the closure
Cartier.  Tensor with the inverse of its restriction at the
identity to rigidify it.  The theorem of the cube extends its
generic cubical structure uniquely; a vertical ambiguity is
pulled back from the base and is removed by rigidification.
The generic polarization and its inverse extend, so the
resulting principal polarization is the same one on every
fibre.  The associated line bundle is relatively ample, by
the abelian-scheme polarization criterion.  Symmetry and all
rigidified generic identifications extend uniquely.  One can
equivalently use the original cubical principal-triple
extension theorem; its hypotheses are satisfied by this
everywhere-good model.

For an ample principal line bundle on an abelian variety,
the higher cohomology vanishes and the space of sections has
dimension one.  Cohomology and base change therefore makes
the relative section module an invertible
\(\mathcal O_{K'}\)-module, compatible with flat base change.
Let \(M\) be the Lie module.  After making all labels
\(\tau:E\hookrightarrow K'\) available, its \(\tau\)-eigenspace
is the kernel of the finite matrix
\[
 M\longrightarrow\bigoplus_a M,\qquad
 v\longmapsto (av-\tau(a)v)_a,
\]
where \(a\) ranges over integral generators of
\(\mathcal O_E\).  This kernel is saturated in the torsion-free
module \(M\), hence is finite projective over the Dedekind
ring.  Its generic rank is the original \(m_\tau=\mathrm{Dw.mult}(\tau)\).
Flatness commutes with the displayed kernel, so this is also
the whole integral eigenspace after localization, completion,
or further number-field extension.  We have not asserted
that the sum of these eigenspaces is the full integral Lie
module at ramified primes; that stronger statement is not
needed and is generally unavailable.
\end{proof}
% END ICL5

% BEGIN ICL6
\paragraph{ICL6. Simultaneous integral frames and their density.}
\label{claim:ICL6}
After one further finite extension, all the projective
modules in ICL5 admit bases.  Starting with these bases and a
theta generator, the union of their integral frame changes
over finite number-field extensions is Zariski dense in the
original generic frame fibre.
\begin{proof}
Over a Dedekind ring a projective module of positive rank is
the sum of free modules and one fractional ideal.  For each
of the finitely many ideals \(I\), choose \(h>0\) and \(a\ne0\)
with \(I^h=(a)\), and adjoin a root \(\alpha^h=a\).  At every
prime of the new number field the valuations of
\(I\mathcal O\) and \((\alpha)\) agree, so the ideals agree.
A compositum principalizes them all.  Rank-zero eigenspaces
need no basis.  The kernel and section-module base-change
statements in ICL5 show that these are bases of the actual
new modules, not of proper finite-index submodules.

For a matrix block of size \(d\geq2\), the integral
elementary matrices \(1+tE_{ab}\), \(t\in\mathbb Z\), are
Zariski dense in each root subgroup: a polynomial with
infinitely many integer zeros is zero.  These subgroups
generate \(\mathrm{SL}_d\).  Diagonal matrices whose entries
are roots of unity supply a Zariski dense subgroup of the
diagonal torus, because a nonzero Laurent polynomial in one
variable has only finitely many roots.  The same argument
handles \(d=1\), and \(d=0\) is trivial.  The Zariski closure
of all these matrices is therefore \(\mathrm{GL}_d\).
Their determinants are algebraic units, so both each matrix
and its inverse are integral.  Products across all blocks,
and roots of unity in the theta-generator scalar, are dense
in the whole product.  Every individual point uses only
finitely many roots of unity and hence is defined over a
finite extension.  Transposing the matrices gives exactly
the source frame-action convention and preserves all these
properties.
\end{proof}
% END ICL6

% BEGIN ICL7
\paragraph{ICL7. A nonzero integral-frame evaluation.}
\label{claim:ICL7}
For every nonzero \(f\in C_w\), there are a number field
\(K_f/K\) and an actual framed point \(y_f\) of \(X_j\) over
\(K_f\) such that
\[
 \operatorname{ev}_R(f,y_f)\ne0,\qquad
 \operatorname{ev}_R(g,y_f)\in\mathcal O_{K_f}\quad(g\in L_w).
                                                               \tag{ICL.4}
\]
The map \(g\mapsto\operatorname{ev}_R(g,y_f)\) is \(K\)-linear.
\begin{proof}
Take the CM point of ICL4 where the regular function defined
by \(f\) on its frame fibre is nonzero.  ICL5--ICL6 give an
integral frame outside its zero locus over some finite
extension; call it \(y_f\).  At every finite prime
\(\mathfrak P\) of \(K_f\), complete its integral model.
The chosen global eigenbases become bases of the entire
eigenmodules over this complete DVR, by the kernel argument
in ICL5.  The chosen theta generator becomes a generator of
the entire section module, by cohomology and flat base change.
Together with the rigidified identifications, these are
literally the witnesses in the source predicate
\(\mathrm{Integral}\).

The point has a \(\mathrm{FromTest}\) presentation: its
number field and finite list of coordinates have a reduced
finite-type presentation over \(K\), and its completion is
obtained from that presentation by base change.  The residue
field of the completion is finite, so
\(\mathrm{IntFr.IsAlg}\) holds.  If
\(\mathfrak p=\mathfrak P\cap\mathcal O_K\), the ordinary
ramification index \(e(\mathfrak P/\mathfrak p)>0\) gives the
source \(\mathrm{OverLa}\) condition.  Thus (ICL.1), applied
to this precise point, is among the tests in (ICL.2).
For \(g\in L_w\) it says that the \(\mathfrak P\)-valuation
of \(\operatorname{ev}_R(g,y_f)\) is nonnegative.  This is
true for every finite \(\mathfrak P\), so the value lies in
\(\mathcal O_{K_f}\).  Naturality identifies the completed
evaluation with the image of the same \(K_f\)-value.
Linearity follows from the source evaluation law.  No
denominator depending on \(g\) has been introduced.
\end{proof}
% END ICL7

% BEGIN ICL8
\paragraph{ICL8. Finite generation at every original weight.}
\label{claim:ICL8}
For every original row, dominant weight, pinned moduli datum
and supplied valid uniformization, the whole lattice
\(L_w\) is finitely generated over \(\mathcal O_K\).
\begin{proof}
The fixed-weight cusp space is finite-dimensional by the
classical analytic finite-dimensionality foundation stated
above.  Its application here uses the original finite set
of arithmetic uniformization components.  The reading map
preserves finite linear independence after extension
\(K\hookrightarrow\mathbb C\).  For each reduced finite-type
test algebra \(A\), a relation between readings makes the
corresponding element of \(A\otimes_{K,\sigma}\mathbb C\)
vanish at every complex point.  This scalar extension is
reduced because \(K\) has characteristic zero, and its
complex points detect its elements, so the element is
zero.  Applying \(K\)-linear coefficient functionals
\(\mathbb C\to K\) to this tensor equality gives a
relation between the original rules on every \(A\).
Independence of those rules makes every coefficient
functional vanish on each complex coefficient, hence
each coefficient is zero.  Thus a finite
independent family of rules maps to an independent family
of analytic cusp forms and is bounded in size by their
finite complex dimension.  This is the literal
\(\mathrm{cuspForms\_finiteDimensional\_of\_goodEmb}\)
application; it uses no rational-theta test family.

Apply ICL2 to the evaluations of ICL7.  This proves finite
generation of \(\mathrm{supLat}(T,j,w)\) itself.  The
argument is valid even if that lattice has not yet been
shown to span and even if the cusp space is zero.  It proves
the first conjunct of the original \(\mathrm{FullAt}\)
without any short-theta spanning or denominator assumption.
\end{proof}
% END ICL8

% BEGIN ICL9
\paragraph{ICL9. Clearing a positive power clears the section.}
\label{claim:ICL9}
Let \(\mathcal R\) be a commutative \(K\)-algebra with the
valuation families of ICL1, and let \(S\subseteq\mathcal R\).
Suppose every \(g\in S\) has a nonzero scalar multiple
nonnegative at every one of these valuations.  If
\(f^q\in\operatorname{span}_K S\) for an integer \(q\geq1\),
then a nonzero scalar multiple of \(f\) is nonnegative at
all of them.
\begin{proof}
Write the given power as a finite linear combination.  For
each term replace its generator by a scalar multiple whose
valuations are nonnegative.  Clear the finitely many new
coefficients with one nonzero \(t\in\mathcal O_K\).
Then \(t f^q\) has nonnegative valuation everywhere.  Fix
one valuation and write \(a=\operatorname{ord}_{\mathfrak p}(t)\geq0\).
If \(v(f)=+\infty\), the assertion for \(tf\) is immediate.
Otherwise \(v(f)=x\in\mathbb R\) and \(a+qx\geq0\).
For \(x\geq0\), also \(a+x\geq0\); for \(x<0\),
\(a+x=(a+qx)-(q-1)x\geq0\).  Thus \(tf\) works
simultaneously for every prime and every frame test.
This is why the clearing scalar is chosen integral.
No extraction of a \(q\)-th root and no enlargement of
the fixed field \(K\) is involved.
\end{proof}
% END ICL9

% BEGIN ICL10
\paragraph{ICL10. Powers reach the precise stable range.}
\label{claim:ICL10}
Let \(f\in C_w\) and \(w.n>0\).  Put
\[
 q=4([E:\mathbb Q]r)+4,\qquad r=D.R.r.
\]
There is a dominant exponent pair \((a',n')\), with
\(n'=q\,w.n\geq4([E:\mathbb Q]r)+4\), such that
\(f^q\) belongs to the original cusp space of that exponent.
\begin{proof}
The exponent pair is the sum of \(q\) copies of the original
one.  The defining weak monotonicity inequalities are
preserved by addition; the evenness condition on the theta
exponent is preserved as well.  Naturality and the Borel
weight law multiply, so the product is an algebraic section
of the summed weight.  On every original complex component,
products of the bounded weighted holomorphic readings are
bounded at the summed weight.  At a rational cusp their
Fourier expansions have the usual positive semidefinite
support, with no boundary constant term for a cusp factor.
In a boundary direction, a sum of these nonnegative
exponents can be zero only when every factor has zero
exponent there.  The cusp factor has no such term.  Thus the
product is again cuspidal.  This is the source
\(\mathrm{CuspMul}\) law, iterated on the original rules;
it does not require a chosen basis of the cusp space.
Finally \(w.n\geq1\) gives the displayed inequality.
\end{proof}
% END ICL10

% BEGIN ICL11
\paragraph{ICL11. The exact identity cover and its pins.}
\label{claim:ICL11}
For the original \((X_j,c,l)\) and an original canonical
connection \(\nabla\) satisfying \(\mathrm{IsNab}\), take
\[
 p_d=1,\quad I=\mathcal O_E,\quad l_2=l,\quad X_2=X_j,
 \quad c_2=c,\quad r_2=1,\quad m_2=1,
 \quad c_\delta=1,\quad\phi_\delta=1.                 \tag{ICL.5}
\]
This is the literal identity \(\mathrm{DeltaDat}\).  It
satisfies \(\mathrm{IsLev}\), \(\mathrm{UnitC}\),
\(\mathrm{PsiTie}(1)\), and \(\mathrm{PhiPin}(1,1,0)\).
Every field-valued evaluation of \(c_\delta\) is one.
\begin{proof}
The divisibility and ideal identities are identities at
the same level, \(\gcd(1,l.M)=1\), and
\(I+\overline I=\mathcal O_E\).  Both marked-section
conditions in \(\mathrm{EtaOK}\) reduce to first powers
and the identity maps.  A framed point lifts to itself;
the line-bundle identifications in \(\mathrm{FrOver}\) are
the original identity rigidifications.  The constant rule
one is a weight-zero section and evaluates to one on every
field point, which has a reduced test presentation.
This verifies the last two \(\mathrm{IsLev}\) clauses,
\(\mathrm{UnitC}\), and the trivial \(\mathrm{PsiTie}\).

Both half quotients \(\mathrm{LaI}\) and
\(\mathrm{LaJ}\) are zero.  Hence there is a single tuple
in the latter, its span is the full zero module, and its
constant function has value one.  In the finite Fourier
formula defining \(\mathrm{PhiPin}\), the exponent of the
primitive root includes \(l.M'/p_d=l.M'\).  Each term is
therefore one independently of the chosen representative
of the zero half quotient.  The sum has one term and is
\(\phi_\delta=1\).  This checks the literal pin, including
the source's representative choices.  The constant rule
has order one and \(3\nmid p_d\).  This construction is
an identity cover; it supplies no higher tuple cover.
\end{proof}
% END ICL11

\subsection{The analytic spanning argument and its exact inputs}

We now state the original inputs still required for this
part.  They are kept visible so that the final lattice
conclusion cannot be mistaken for an unconditional closure
of the retained RTG producer.

The principal-chart input supplies, on the same \(X_j,c\),
a valid uniformization \(U'\) whose components have the
same Hermitian matrix and imaginary scalar, whose lattice
stabilizers contain the original principal level, and
whose scalar multiplier \(\chi\) is one.  It also identifies
its metric with the original one.  This is the precise
\texttt{SupFull.exists\_unf} construction and metric
comparison used in the source; it is not implied by an
arbitrary \(\mathrm{IsUnf}\).

The level-system input supplies the actual lattices
\(\mathcal L\), unitary group \(\Delta\), vacuum functions
\(\Phi(L,\Lambda)\), and unit multipliers \(\mu(\gamma,L)\)
with all fields of the original \(\mathrm{IsSystem}\).
In particular the finite class sum uses precisely its
positive weights \(\mathrm{wt}(M,\Lambda)\), and its norm
field has a summable off-diagonal majorant of total
strictly less than one.  The numerical range is
\(n\geq4([E:\mathbb Q]r)+4\), with \(n\) even.

The raw Gaussian input consists of the source
\texttt{kernel\_identity}, the integrability of that
kernel on a sufficiently thin orthogonal slab, and its
reproducing formula there.  With an invariant finite
harmonic family \(P_b\), write its kernel as \(\mathcal K\).
The formula and normalization used below are
\begin{align}
 \kappa\sum_{\gamma\in\Delta}
 C(L,\gamma M)\overline{\mu(\gamma,M)}
 \mathcal K(x,\gamma y)
 &=\sum_{[\Lambda],b}\mathrm{wt}(M_0,\Lambda)
       \Theta_{\Lambda b}(L,x)
       \overline{\Theta_{\Lambda b}(M,y)},             \tag{ICL.6}\\
 \int_S\mathcal K(x,y)H(y)\,dy&=\rho H(x),\qquad
                    \kappa\ne0,\quad\rho>0.          \tag{ICL.7}
\end{align}
Here \(M_0\) denotes the integer level, to distinguish it
from a lattice \(M\).  The series in (ICL.6) has its
asserted sum, and the kernel is absolutely integrable on
\(S\).  All measures, Gram matrices, and oscillator
normalizations in these formulas are the original ones.
The following claims prove the passage from these raw
inputs to spanning; they do not assume a spanning theorem.

% BEGIN ICL12
\paragraph{ICL12. The class-sum operator has a bounded inverse.}
\label{claim:ICL12}
For the actual supplied level system, set
\[
 C(L,M)=\sum_{[\Lambda]}\mathrm{wt}(M_0,\Lambda)
       \langle\Phi(L,\Lambda),\Phi(M,\Lambda)\rangle.
                                                               \tag{ICL.8}
\]
There are \(a>0\) and \(0\leq\theta<1\) such that
\[
 C(L,L)=a,\qquad
 \sum_{M\ne L}|C(L,M)|\leq\theta a.                  \tag{ICL.9}
\]
On the space of original bounded holomorphic forms with
the prescribed Borel weight and \(\mu\)-covariance,
\((\mathcal CH)(L,x)=\sum_M C(L,M)H(M,x)\)
is invertible and its inverse has norm at most
\(a^{-1}(1-\theta)^{-1}\) in the weighted supremum norm.
\begin{proof}
The class set is finite and nonempty, its weights are
positive, and their sum \(m_0\) is positive.  If the
vacuum norm in \(\mathrm{IsSystem.norm}\) is \(a_0>0\),
the diagonal is \(a=m_0a_0\).  Its majorant \(b(L,M)\)
gives \(|C(L,M)|\leq a b(L,M)\).  Summing gives
(ICL.9), including absolute convergence.  If
\(\mathcal L\) is empty the whole assertion about forms
is the assertion about the zero space; otherwise the
nonnegative off-diagonal sum also permits \(\theta\geq0\).
The unitary invariance of the vacuum inner product gives
\[
 C(\gamma L,\gamma M)\mu(\gamma,L)
                \overline{\mu(\gamma,M)}=C(L,M).
\]
Together with \(|\mu|=1\), this proves that
\(\mathcal C\) preserves the exact covariance convention
\(H(\gamma L,\gamma x)=\mu(\gamma,L)^{-1}H(L,x)\).

The space of bounded forms is complete in the weighted
supremum norm.  A Cauchy sequence converges uniformly after
weighting; on every compact subset of the frame domain
the weight and its inverse are bounded, so the limit is
locally uniformly holomorphic and satisfies the same
transformation laws.  Write \(\mathcal C=a(1+Q)\).
(ICL.9) implies \(\|Q\|\leq\theta<1\), and the uniformly
convergent series
\(a^{-1}\sum_{s\geq0}(-Q)^s\) is its two-sided inverse.
In particular the inverse preserves the holomorphic and
covariance conditions.  This gives the original Neumann
step with its actual class-sum coefficients.
\end{proof}
% END ICL12

% BEGIN ICL13
\paragraph{ICL13. Unfolding a finite-rank kernel.}
\label{claim:ICL13}
Let a countable group \(\Gamma\) act on a countable set
\(I\) and preserve a measure on \(Y\).  Suppose that the
diagonal action is free almost everywhere and has measurable
fundamental pieces \(D_i\) whose translates partition
\(I\times Y\) almost everywhere.  Let
\(f_z(i,y)\) satisfy
\(\sum_i\int_Y|f_z(i,y)|\,dy<\infty\) for every
parameter \(z\) in a set \(Z\).  If on the fundamental
pieces its orbit sum has finite rank in \(z\),
\[
 \sum_{\gamma\in\Gamma}f_z(\gamma i,\gamma y)
              =\sum_{b=1}^s e_b(z)t_b(i,y),            \tag{ICL.10}
\]
then \(z\mapsto\sum_i\int_Y f_z(i,y)\,dy\) lies in
the finite linear span of the functions \(e_b\).
For the original lattice/frame action the required
fundamental pieces exist and the measure is the original
Lebesgue measure on frames, restricted to the slab.
\begin{proof}
Absolute summability permits Tonelli and the changes of
variables \((i,y)\mapsto(\gamma i,\gamma y)\).  The
fundamental pieces partition the orbit space, so the
integral of the orbit sum on these pieces equals
\(\sum_i\int_Y f_z\).  There is a small integrability
point here: the individual \(t_b\) need not initially be
integrable.  Replace the \(e_b|_Z\) by a basis of their
span, and choose finitely many evaluation points
\(z_1,\ldots,z_t\) for which the evaluation matrix is
invertible.  Solving that matrix system expresses each
corresponding coefficient function on the pieces as a
fixed linear combination of the integrable orbit sums
at \(z_i\).  Each coefficient is consequently integrable.
Integrating (ICL.10) now proves the assertion.

For the lattice/frame application, give the countable
set of full lattices the discrete topology.  The action
on its product with the invertible-frame domain is free:
an element fixing a full frame is the identity.  It is
properly discontinuous.  On a compact set there are only
finitely many lattice labels; a matrix carrying one
fixed lattice to another has discrete coefficients in
fixed lattice bases, and the frame and inverse-frame
bounds bound these coefficients.  Only finitely many
such matrices can occur.  Choose a countable collection
of small open sets with disjoint translates and remove
earlier translates to obtain Borel fundamental pieces.
Unitary left multiplication has real Jacobian
\(\prod_\tau|\det\gamma_\tau|^{2r}=1\), and preserves
the slab defined by the original Gram matrices.  Thus
these are precisely the measure-preserving changes of
variables required above, with no rescaling of measure.
\end{proof}
% END ICL13

% BEGIN ICL14
\paragraph{ICL14. The raw original kernel implies theta spanning.}
\label{claim:ICL14}
Assume the level-system and raw Gaussian inputs just
stated.  Every original bounded holomorphic form \(F\)
is a finite linear combination of the actual functions
\(\Theta_{\Lambda b}\), simultaneously for every
\(L\in\mathcal L\) and every frame in the domain.
\begin{proof}
Set \(H=\mathcal C^{-1}F\) using ICL12.  Choose a slab
small enough for (ICL.6)--(ICL.7) and for the source
uniform Borel-factor bounds.  These bounds and boundedness
in the weighted norm give one ordinary bound for
\(|H(L,y)|\) on the slab, independent of \(L\).
For \(z=(L,x)\), apply ICL13 to
\[
 f_z(M,y)=C(L,M)\mathcal K(x,y)H(M,y).
\]
The diagonal term is integrable by the kernel bound;
the off-diagonal terms are dominated by its integral
times the uniform bound for \(H\) and the summable
series in (ICL.9).  This verifies the absolute
summability hypothesis, rather than exchanging
uncontrolled infinite sums.

In the orbit sum, the covariance of \(H\) contributes
\(\mu(\gamma,M)^{-1}=\overline{\mu(\gamma,M)}\).
Therefore (ICL.6) makes that sum a finite linear
combination, in \(z\), of
\(\Theta_{\Lambda b}(L,x)\).  Its coefficient at
\((M,y)\) is exactly
\(\kappa^{-1}\mathrm{wt}(M_0,\Lambda)
\overline{\Theta_{\Lambda b}(M,y)}H(M,y)\).
ICL13 now gives a finite expansion of the unfolded
integral.  On the other hand (ICL.7) and absolute
summability evaluate that integral as
\[
 \rho\sum_M C(L,M)H(M,x)=\rho F(L,x).
\]
Divide by the nonzero real number \(\rho\).  The result
is the asserted expansion, with coefficients independent
of both \(L\) and \(x\).  Its class representatives,
weights, phases and measure are those of (ICL.6).
This proves the source analytic spanning implication
from raw kernel and level-system data.  It does not
prove those raw data from the common geometric family.
\end{proof}
% END ICL14

% BEGIN ICL15
\paragraph{ICL15. Integral harmonic polynomials span the actual space.}
\label{claim:ICL15}
For any full original lattice \(\Lambda\subset E^n\),
the polynomials in \(\mathrm{harmL}(c,\Lambda,a)\)
span the original harmonic highest-weight space after
extension from \(K\) to its fixed complex embedding.
\begin{proof}
All polynomials under consideration have one fixed
finite total degree.  The row-weight, highest-weight
and harmonic conditions are kernels of finitely many
linear maps over \(K\): the original field \(c.F\)
contains the prescribed embeddings of \(E\) and their
conjugates.  Kernel formation commutes with a field
extension.  Choose a \(K\)-basis of that finite
solution space.

Choose a positive integer \(N\) with
\(N\mathcal O_E^n\subseteq\Lambda\), possible because
\(\Lambda\) is full.  Each variable multiplied by
\(N\) is one of the lattice-linear generators of
\(\mathrm{symO}(c,\Lambda)\).  For a basis polynomial
of degree \(d\), multiplication by \(N^d\), followed
by an integral scalar clearing its finitely many
coefficients, puts it in this subring.  The multiplier
is nonzero and preserves all the linear equations.
The resulting scaled polynomials belong to the
literal \(\mathrm{harmL}\) and retain their span.
Renaming their variables by the original letter
dictionary gives the corresponding row-highest
polynomials on the analytic side.  No independently
chosen harmonic representation is substituted.
\end{proof}
% END ICL15

% BEGIN ICL16
\paragraph{ICL16. The exact short-series span over the original field.}
\label{claim:ICL16}
Retain the original principal-chart and finite-level
dictionary input on \(U'\), the identity cover (ICL.5),
the chart-independent construction of its short rules,
and the raw inputs of ICL14.  If
\(n\geq4([E:\mathbb Q]r)+4\), then every section
\(f\in\mathrm{sec}(X_j,c,a,n)\) whose readings in
\(U'\) are bounded belongs to
\[
 \operatorname{span}_K\mathrm{shortGen}
                     (X_j,c,\nabla,1,0,a,n,d).        \tag{ICL.11}
\]
\begin{proof}
Here is the full use of the dictionary input.  It supplies
one reference Hermitian space, original component
lattices \(L_k\) in distinct \(\Delta\)-orbits, and
unitary coordinate changes \(\gamma_k\).  Its vacuum
multiplier is one on every stabilizer of an \(L_k\).
Pull the \(k\)-th reading back by \(\gamma_k\) and
divide by its fixed nonzero comparison scalar \(C_k\).
The principal-level inclusion, the trivial scalar
character, and the stabilizer multiplier show that it
is invariant under the stabilizer with exactly the
given Borel weight.  Extend it to the orbit of \(L_k\)
by the prescribed \(\mu^{-1}\)-covariance, and set it
zero on the other orbits.  The cocycle identity and
stabilizer invariance prove well-definedness.  There
are finitely many \(k\), so the resulting form has
a uniform weighted bound.  ICL14 expands it into
finitely many theta terms.

The lower-row and lattice-translation part of the
retained original dictionary rewrites those terms
using the actual allowed lattices \(\Lambda\) and
the original level \(M=l.M\).  ICL15 expands their
harmonic polynomials using the literal integral
\(\mathrm{harmL}\).  The retained short-rule
construction supplies, for each resulting \(\Lambda,N,P\),
a section \(g\) in \(\mathrm{shortGen}\).  Its exact
reading identity is
\[
 \operatorname{rd}_{U'}(g)_k(x)
   =C_k u_{\Lambda,N}
       \Theta_{\Lambda,P}(L_k,\gamma_k x),
            \qquad u_{\Lambda,N}\ne0.                 \tag{ICL.12}
\]
The scalar \(u_{\Lambda,N}\) is independent of the
component \(k\).  Dividing each theta coefficient by
that scalar gives one finite complex linear
combination of readings of original short rules,
valid on every component.  Off the frame domain all
these readings and the given reading are zero by
their definitions, so equality holds on the entire
source function space.

It remains to descend the coefficients; mere
injectivity of the reading would not suffice.  Its
preservation of finite \(K\)-linear independence
after complex extension was proved in ICL8 directly
from naturality and reduced test rings.  If (ICL.11)
failed, choose a basis of the span of the finitely
many rules in the complex expression and append
\(f\).  This is \(K\)-independent, while its readings
are complex-dependent, a contradiction.  Therefore
the coefficients can be chosen in the original
\(K\).  All assertions here concern the identity
cover; no higher \(\mathrm{CoversAt}\), arbitrary
tuple class sum, or rescaled phase has been supplied.
\end{proof}
% END ICL16

\subsection{The finite-place calculation and the full-lattice endpoint}

The remaining integral-theta input is the following original,
uniformly quantified statement.  For every allowed short
rule \(g=g_{\Lambda,N,P,x_0}\), and every test in (ICL.1),
one may pass to a finite extension of the complete DVR,
of ramification degree \(b>0\), and use the original
higher-level integral chart and torsion data.  With
\(e'=be\), its normalized valuation \(\nu\) satisfies
\(\nu(\iota' z)=e^{-1}\operatorname{ord}_{V}(z)\).
The precise short-value formula, after grouping the
original torsion sum, is
\[
 \iota'(\beta_H g(y))
       =\zeta\,c_0^{-1}\sum_{q}w_q X_q,
 \quad c_0=\#\bigl((\mathrm{Fin}\,n)\to
                              \mathrm{LevQ}(\mathrm{Dw},l_2)\bigr),
                                                               \tag{ICL.13}
\]
where
\[
 X_q=\iota_2\bigl(\mathrm{mIdx}(\Lambda,N,r)^{-1}
                                     N^{-\deg P}\bigr)
       \sum_i\mathrm{hv}\bigl(\mathrm{sqN}(F_{qi}),P\bigr).
\]
The scalar \(\beta_H=\mathrm{betaH}(c,\Lambda)\ne0\)
is the literal source scalar, transported by \(d.eF\)
on the cover.  The weights \(w_q\) are its literal finite
Fourier weights and are integral.  The scalar \(\zeta\)
has a positive finite order.  The torsion points satisfy
the source \(\mathrm{InZh}\), geometric injectivity,
\(\mathrm{IsUUOver}\), \(\mathrm{AllRat}\), and
half-ideal killing conditions, not merely an abstract
finite-set model.

For a squarefree \(M=l_2.M\), put \(\alpha=1\) if
the residue characteristic \(\ell\) does not divide
\(M\), and \(\alpha=\ell\) otherwise.  The local
jet input gives, for every \(q\),
\begin{align}
 \nu\bigl(\operatorname{nrm}(\alpha)X_q\bigr)&\geq-B,
       &\operatorname{nrm}(\alpha)&=\alpha^{d_w},       \tag{ICL.14}\\
 B&=\frac{d_w}{2}\sum_{p\mid\mathrm{condN}}
                        \nu(p)\,\mathrm{vdH}(p),
       &d_w&=\mathrm{DF}(\mathrm{Dw},w).               \tag{ICL.15}
\end{align}
This is the exact combination of the source
\texttt{class\_bound} in \texttt{ExistsSmulShortNonneg}
and \texttt{HostChart.chBound\_host}.  It includes the
ramified case and every integral frame.  It is retained
as an original jet producer, not proved here by an
appeal to integrality at good primes.  The proof below
shows exactly why this bound, with these normalizations,
is sufficient and gives a single global clearing scalar.

% BEGIN ICL17
\paragraph{ICL17. A normalized finite average has a uniform bound.}
\label{claim:ICL17}
Let \(\nu\) be a normalized discrete valuation, let
\(u,c\ne0\), let \(w_q\) be integral, and let
\(\zeta^m=1\) with \(m>0\).  If
\(\nu(uX_q)\geq-B\) for all members of a finite set,
then
\[
 \nu\left(\zeta c^{-1}\sum_qw_qX_q\right)
                 \geq-B-\nu(u)-\nu(c).                \tag{ICL.16}
\]
The assertion remains valid when the average is zero.
\begin{proof}
The root-of-unity equation implies \(\zeta\ne0\) and
\(m\nu(\zeta)=0\), hence \(\nu(\zeta)=0\).
Each \(w_q\) has nonnegative valuation because it is
integral over \(\mathbb Z\) and the valuation ring is
integrally closed.  Thus
\(\nu(w_quX_q)\geq-B\), and repeated use of the
ultrametric inequality gives the same bound for their
sum.  Multiplication by \(c^{-1}u^{-1}\) subtracts
\(\nu(c)+\nu(u)\).  If any summand or the full sum
is zero its valuation is \(+\infty\), which is
compatible with each inequality.  No archimedean
triangle estimate or factor equal to the cardinality
of the summation set is inserted; the only cardinality
denominator is the literal \(c\) already in the formula.
\end{proof}
% END ICL17

% BEGIN ICL18
\paragraph{ICL18. One explicit scalar clears every original place.}
\label{claim:ICL18}
Assume (ICL.13)--(ICL.15) for a permitted original cover
with squarefree \(M=l_2.M\), and use the original host
bound \(0\leq\mathrm{vdH}(p)\leq\Delta_E\), where
\(\Delta_E=|\operatorname{disc}(E)|\).  Define
\[
 P_0=\prod_{p\mid\mathrm{condN}}p,\qquad
 A_w=P_0^{\Delta_E d_w}M^{d_w}c_0,\qquad
 s=A_w\beta_H\in K.                                 \tag{ICL.17}
\]
Then \(s\ne0\) and \(v(sg)\geq0\) for every
original \(v\in\mathrm{valSet}(\mathfrak p)\), at
every \(\mathfrak p\).  Thus the exact short-generator
denominator conclusion holds with the scalar (ICL.17).
\begin{proof}
The prime product is nonzero, \(M\ne0\), the finite
level quotient is nonempty and has positive cardinality,
and \(\beta_H\ne0\) by its source definition.  Hence
\(s\ne0\).  The stated host bound follows directly
from its definition
\(\mathrm{vdH}(p)=\operatorname{ord}_p(\Delta_E)/[E:\mathbb Q]\):
the denominator is at least one, and for each prime
\(p\) used here, \(p^{\operatorname{ord}_p(\Delta_E)}\)
divides the positive integer \(\Delta_E\).  Thus
\(0\leq\operatorname{ord}_p(\Delta_E)\leq\Delta_E\),
which proves both inequalities without a new local
estimate.  Fix any original test, make the finite
extension supplied by (ICL.13), and retain the normalized
valuation \(\nu=\operatorname{ord}_{V'}/e'\).
Because \(M\) is squarefree, \(\ell^2\nmid M\);
the chosen \(\alpha\) is a positive integer dividing
\(M\).  Write \(M=\alpha\alpha'\).
Apply ICL17 to (ICL.13), with
\(u=\operatorname{nrm}(\alpha)=\alpha^{d_w}\)
and \(c=c_0\).  It gives
\[
 \nu(\beta_Hg(y))\geq-B-d_w\nu(\alpha)-\nu(c_0).
\]
Add
\(\nu(A_w)=\Delta_Ed_w\sum_{p\mid\mathrm{condN}}\nu(p)
                  +d_w\nu(M)+\nu(c_0)\).
The resulting lower bound is
\[
 d_w\sum_{p\mid\mathrm{condN}}
        \nu(p)\left(\Delta_E-\frac{\mathrm{vdH}(p)}2\right)
                      +d_w\nu(\alpha')\ \geq\ 0.      \tag{ICL.18}
\]
Every factor is nonnegative.  Natural integers are
integral in the valuation ring, and the host bound
implies the displayed coefficient is nonnegative.
Here the literal degree is
\(d_w=\sum_\tau\sum_{i<\mathrm{Dw.mult}(\tau)}w(\tau,i)\in\mathbb N\);
its nonnegativity is part of its definition, including
when that sum is zero.
Normalized valuation is preserved by the finite
extension, so the original test also gives a
nonnegative value.  The scalar (ICL.17) was fixed
before the test was chosen.  It is therefore one
scalar for all places and all original frames.

For the identity cover of ICL11, all of its Fourier
weights are powers of the original \(\zeta_{l.M}\)
and are integral, \(c_\delta=1\), and one can take
the root order to be one.  Its level remains
\(l.M=g_0g_1\), squarefree by the two distinct
original pin primes.  Hence the assumptions involving
level, phase and squarefreeness used here are exactly
those needed in the source application.
\end{proof}
% END ICL18

In the next three claims the retained producer inputs are
quantified over every original row and every dominant
exponent used in the argument, including each positive
multiple introduced in ICL10.  The original \(D,T,c\),
field embeddings, metric, measure and phases are fixed
throughout.  A chart or short-rule witness is allowed to
depend on the row and exponent precisely as in the
displayed original producer statements; its existence
is not asserted to follow from the finite-generation
half of this chapter.

% BEGIN ICL19
\paragraph{ICL19. The original positive-weight full lattice, with its exact remaining inputs.}
\label{claim:ICL19}
Assume \(T.Q.\mathrm{Pins}\), a valid original
uniformization at the distinguished good embedding,
and the principal-chart, level-system, Gaussian,
short-rule/dictionary and integral-jet inputs specified
above.  For every original row \(j\), every dominant
weight \(w\) with \(w.n>0\),
\[
 \mathrm{FullAt}(T,j,w)
 \quad\text{and}\quad
 \mathrm{FullLat}(C_w,\mathrm{valSet},\mathrm{isVal})
                                                               \tag{ICL.19}
\]
hold on the original \(K\)-space and its whole
valuation-nonnegative lattice.
\begin{proof}
Finite generation is ICL8.  To prove spanning, fix
\(f\in C_w\), and take \(q,a',n'\) from ICL10.
The original cusp readings are bounded.  Their
boundedness passes to the supplied principal charts
without a change of metric.  More explicitly, at
orthonormal frames the theta metric is one and the
\(\tau\)-Lie metric is
\(2\pi|\tau(\xi_0)|\) times the identity.  A Borel
change between the two presentations has section
scalar of modulus one; its diagonal Lie scalars
have moduli equal to the square roots of the ratios
of these fixed positive constants.  The prescribed
weight multiplies the old bounds by the corresponding
finite product.  Taking a sum of these bounds over
the finitely many old components gives a bound on
each new component.  This is the exact metric
comparison used by \texttt{isBddSec\_of\_isUnf}.

ICL16 now writes \(f^q\) in the \(K\)-span of
the original short generators at \((a',n')\).
ICL18 supplies a nonzero clearing scalar for each
one.  ICL9 gives \(t\ne0\) with all valuations
of \(tf\) nonnegative.  Since \(C_w\) is a
\(K\)-subspace, \(tf\in L_w\), and
\(f=t^{-1}(tf)\) belongs to \(K L_w\).  The
reverse inclusion follows from \(L_w\subseteq C_w\).
This proves \(\mathrm{FullAt}\).  The source
\(\mathrm{FullLat}\) places the same lattice inside
the subtype \(C_w\); the inclusion is an isomorphism
onto \(L_w\).  Finite generation and spanning
transport across that inclusion by ICL1.  This
proves both statements in (ICL.19), with no added
flatness, nonvanishing, or integral-basis premise.
\end{proof}
% END ICL19

% BEGIN ICL20
\paragraph{ICL20. The literal \texttt{LatFull} and canonical-weight specializations.}
\label{claim:ICL20}
Under the inputs of ICL19, let \(T:\mathrm{CoreJ}(D)\)
and \(\kappa\) satisfy the original \(\mathrm{LawsJ}\)
and \(\mathrm{KapIs}\).  Then the literal predicate
\(\mathrm{LatFull}(T,\kappa)\) holds.  The same proof
gives the canonical-weight
\texttt{IntegralCuspFormsFullLattice} endpoint.
\begin{proof}
Fix an arbitrary \(i\in\mathrm{Fin}(D.R.m)\) and
put \(j=i+1\).  The rank inequalities in \(D\)
give \(D.R.nV>0\).  Choose a member of its
nonempty CM embedding family.  The corresponding
\(\mathrm{KapIs}\) equation gives
\((\kappa_i).n=D.R.nV\); the weight is therefore
positive.  The fields of \(\mathrm{LawsJ}\)
give the original pins and valid uniformization
at this row and distinguished embedding.  ICL19
applies and says exactly
\[
 \mathrm{FullLat}\bigl(T.\mathrm{cuspA}(i+1,
                 \mathrm{wtA}(\kappa_i),(\kappa_i).n),
                 \mathrm{valSet}(X_{i+1},c,\mathrm{algCls})\bigr).
\]
The row was arbitrary, so this is the definition
of \(\mathrm{LatFull}\).  If there are no rows the
universal assertion is vacuous.  For the canonical
weight \(\mathrm{kaWt}\), use its literal equation
\(\mathrm{kaWt}.n=D.R.nV\) and its original
\(\mathrm{kappaExp}\).  The lattice is the same
\(\mathrm{supLat}\) by definition.  This supplies
that specialization without altering the common
coefficient record or selecting a new weight family.
The original geometric inputs named before ICL19
remain premises of both specializations.
\end{proof}
% END ICL20

% BEGIN ICL21
\paragraph{ICL21. Exact transport to the original \(M_K\) and both Boolean levels.}
\label{claim:ICL21}
Retain all inputs of ICL20, the original
\(t:\mathrm{FWTie}(T,R,\kappa,v_0)\), and an
arbitrary original face \(F\).  Then
\((R.FW_i).\mathrm{MKFull}\) holds for every row,
and \((R.\mathrm{toCx}(F)).\mathrm{MLat}(b)\)
holds for \emph{each} \(b\in\{\mathrm{false},
\mathrm{true}\}\).
\begin{proof}
Use the exact map \(\iota_i\) supplied by the tie.
Its rational range is the original \(W_{K,i}\).
The point and norm parts of that same tie identify
\(\iota_i(f)\in M_{K,i}\) with nonnegativity of
every original valuation of \(f\).  This is the
full test equivalence in RTG22, not just an
inclusion of one integral image.  Hence
\[
 M_{K,i}=\iota_i(L_{\kappa_i}).                         \tag{ICL.20}
\]
ICL20 and the source map now give finite generation
and the full rational span of this module.

For either Boolean level, the original good/type
subspace and coordinate map define \(M_j\) and
\(M_{\mathrm{coord}}\) by intersection and exact
inverse image.  RTG23 transports a full ambient
lattice to these original subspaces and through
the specified coordinate isomorphism.  Apply it
using (ICL.20), once with \(b=\mathrm{false}\)
and once with \(b=\mathrm{true}\).  The same
\(T,R,\kappa,v_0,F\), field embedding, group
operators and phase data are used in both
applications.  Thus this is the literal
\texttt{mkFull\_of\_tie}/\texttt{mLat\_of\_tie}
consequence.  It changes neither the projector
nor either level's coefficientwise map.
\end{proof}
% END ICL21

\subsection{What this packet supplies and what it still requires}

ICL8 supplies the actual finite-generation half of
the original geometric lattice, from the stated
classical CM, abelian-scheme and fixed-weight
automorphic foundations.  It needs none of the
original short-theta producers.  Its simplification
does not claim the source's stronger quantitative
\(\mathrm{CMPt.Good}\) or bounded-radius theorem.
ICL11 constructs the actual identity cover.  ICL12--ICL18
prove the analytic unfolding, coefficient descent,
and global-denominator implications at the original
normalizations.  ICL19--ICL21 are complete implications
from the explicitly retained producer inputs; those
inputs have not all been supplied in this packet.

In particular, the actual principal-chart construction,
the common level system, the raw Gaussian identity and
reproducing estimates, the short-rule gluing and its
simultaneous reading dictionary, and the ramified
integral jet estimate remain a quantified frontier
for the spanning half.  Their exact source ranges
are in the contract and foundation manifests.
Consequently the whole original \texttt{LatFull}
supplier is not marked unconditionally reconstructed.

The separate higher \(\mathrm{CoversAt}\), all-tuple
\(\mathrm{ClsEx}\) with phase rescaling, one
\(\mathrm{PointIdentCore}\) chosen before all tuple
records, and base \(\mathrm{ClassSumsSpan}\) also
remain.  This packet closes neither
\texttt{DESIGN\_CHOICES\_COMMON} nor
\texttt{BASE\_THETA\_RATIONAL\_TESTS}, and it
does not infer the outer common-family quantifiers,
eta agreement, or the stronger second division law.

\paragraph{Primary literature used for the classical inputs.}
J.-P. Serre and J. Tate, \emph{Good reduction of abelian
varieties}, Annals of Mathematics 88 (1968), 492--517,
especially Theorem 6(a),
\url{https://annals.math.princeton.edu/1968/88-3/p05}.
For the CM algebraicity and uniformization formulation,
J. S. Milne, \emph{Complex Multiplication}, Chapters I--II,
\url{https://www.jmilne.org/math/CourseNotes/CM.pdf}.
For the cubical extension and relative section-module
framework, L. Moret-Bailly, \emph{Pinceaux de varietes
abeliennes}, Asterisque 129 (1985), Chapters I and VI,
\url{https://www.numdam.org/item/AST_1985__129__1_0/}.
For fixed-type automorphic finite-dimensionality,
Harish-Chandra, \emph{Automorphic forms on semisimple
Lie groups}, Lecture Notes in Mathematics 62 (1968),
Theorem 1,
\url{https://doi.org/10.1007/BFb0098434}; the reductive
form is also stated and proved in A. Borel,
\emph{Automorphic forms on reductive groups}, Section 8.4,
\url{https://homepages.math.uic.edu/~rtakloo/papers/borel/borel3.pdf}.
The manuscript retains these foundational results;
the new proofs above spell out their application to
the exact integral-frame tests rather than claiming
to reprove their general theorems.
